\section{Процесс \code{Retrieve}}\label{PROCESSES.Retrieve}

Процесс~\code{Retrieve} выполняют отправитель не до конца обработанного запроса
на выпуск сертификата и УЦ.  Отправитель~--- это сторона, которая выполнила один
из процессов~\code{Enroll}, \code{Reenroll} или~\code{Spawn} и по его окончании
получила статус~\code{waiting}. Сохранив идентификатор~\code{requestId} своего
запроса, отправитель с помощью~\code{Retrieve} может завершить его обработку.

Процесс~\code{Retrieve} состоит из следующих процедур:
\begin{itemize}
\item
отправка запроса УЦ;
\item
обработка запроса;
\item
возврат ответа;
\item
обработка ответа.
\end{itemize}

Запрос УЦ представляет собой контейнер~\code{BPKIRetrieveReq}, который определен
в~\ref{FMT.BPKIRetrieveReq}. Отправитель записывает в компонент~\code{requestId}
контейнера идентификатор первоначального запроса, а в компонент~\code{nonce}~---
случайную строку октетов. Контейнер отправляется УЦ.

УЦ проверяет наличие присланного в контейнере идентификатора~\code{requestId} в
списке сохраненных пар <<идентификатор запроса~--- сертификат отправителя>>.
Этот список формируется при обработке запросов в процессах~\code{Enroll},
\code{Reenroll} и~\code{Spawn}.

Если идентификатор не найден, то УЦ возвращает контейнер~\code{BPKIResp} со
статусом~\code{rejection}.
%
Если идентификатор найден, но выпуск сертификата по соответствующему запросу все
еще не завершен, то в контейнере возвращается статус~\code{waiting}.
%
Если идентификатор найден, но выпуск сертификата завершен с ошибкой, то в
контейнере снова возвращается статус~\code{rejection}.
%
Во всех случаях~\code{requestId} и~\code{nonce} из запроса переносятся в
соответствующие компоненты~\code{BPKIResp}.

Ответ~\code{BPKIResp} подписывается агентом УЦ. Подписанный ответ оформляется
как контейнер~\code{SignedData}.

Если сертификат все-таки выпущен, то УЦ конвертует его на открытом ключе из
сертификата отправителя, соответствующего~\code{requestId}. Ответ отправляется в
виде контейнера~\code{EnvelopedData}.

Отправитель обрабатывает ответ УЦ так, как если бы он был получен в том
процессе, в котором был отправлен первоначальный запрос. Дополнительно
отправитель проверяет компонент~\code{nonce} в ответах типа~\code{BPKIResp}~---
его значение должно совпадать с первоначальной синхропосылкой, выбранной
отправителем.

